% ---------------------------------------------------------------------------
\section{Imagination that can be trusted}\label{sec:imag}
% ---------------------------------------------------------------------------
Prediction structure and imagination are the natural way to give a selector a sense of consequence: rather than asking ``which candidate resembles what the expert did'', ask ``what happens if I do this''.
The obstacle is identification. A world model trained on expert logs sees each action only in the states where the expert chose it, so its answer for any other action is an extrapolation.
Our programme measured two consequences of trusting such extrapolations: a rollout can be perfectly self-consistent for a candidate that is implausible, and rollout confidence can rise while rollout fidelity falls\evI.
The imagination stream of this design (repository manifest in \S\ref{sec:prelim}) added a third, uncomfortable finding that shapes the whole section: in a controlled toy, \emph{most of what looks like the benefit of imagination is the shape of the training target}, which costs nothing at inference. HiCAP-I is therefore built in three tiers, cheapest and safest first.

\paragraph{Tier 0 (training only): outcome-aware soft targets.}
The expert's one-hot label is the worst target the programme has measured for a selector\evI. Tier 0 replaces it by a listwise target built from per-candidate outcomes, which labels may use freely (ego and other agents),
\begin{equation}
y(a)\;=\;0.7\,\frac{\exp(-J_{\text{teacher}}(a)/T)}{\sum_{a'}\exp(-J_{\text{teacher}}(a')/T)}\;+\;0.3\,\mathbb 1[a=a^\star],
\end{equation}
where $J_{\text{teacher}}$ is a cost computed from the logged future of other agents and the candidate's known kinematics, with one head per metric family (safety, comfort, progress, lateral). It has zero inference cost. In the toy of \S\ref{sec:prelim} it cut regret from $7.00$ to $3.43$ (paired difference $3.57$, 95\% interval $[2.99,4.19]$), more than any imagination arm added on top.

\paragraph{Tier 1 (every tick): imagine the exogenous part, integrate the ego part analytically.}
For a candidate $a$ the outcome is $(W,e_a)$: the future scene $W$ and the ego path $e_a=\kappa(a,e_t)$, a \emph{known} function of the entry and the ego state. Only $W$ needs to be learned, and under the non-reactive assumption of Proposition~\ref{prop:ident} it does not depend on $a$, so it is identified by logs alone. The exogenous part is represented by a handful of descriptors, not by embeddings or frames (a future frame contains the ego's own motion and would smuggle the action into the target):
an E-head $\mathbb R^{512}\!\to\!\mathbb R^{1024}\!\to\!\mathbb R^{24}$ reading the state only outputs the lead gap, closing speed and lead acceleration at $0.5,1,2,3$\,s, a posterior over three to five behaviour modes (for instance the lead brakes, holds, accelerates) and onset-time quantiles. The cost of a candidate is the \emph{expectation over the mode and the timing},
\begin{equation}
J_1(s,a)=\mathbb E_{m,\tau\sim\hat p(\cdot\mid s)}\big[\,c\big(e_a,\ x_{m,\tau}\big)\big],
\end{equation}
with $c$ an analytic time-to-collision and clearance cost evaluated on the unicycle integration of the candidate. Marginalising \emph{every} hidden factor matters: a sharp mode posterior with a mean timing was over-confident in the toy ($+0.7$ regret). Tier 1 needs logs only, no counterfactual data, and it is the one imagination arm that survives their absence (regret $1.80$ against $47$--$72$ for the unstructured arms). Its cost is $\approx3$\,MFLOP per tick on the pruned set.

\paragraph{Tier 2 (gated, top-8, only where counterfactual supervision exists).}
Two cost heads on $\phi(s,a)=\operatorname{norm}(\mathbf s+R\,\mathbf v(a))$ regress per-family \emph{counterfactual} costs of a candidate, obtained from a rule teacher or a simulator (Proposition~\ref{prop:future} shows why the score cannot be an outcome-InfoNCE: it is identified only up to a per-candidate constant and cannot rank candidates, and a candidate-normalised softmax needs counterfactual data). Outcome InfoNCE is kept only as auxiliary representation shaping, and its targets are the exogenous descriptors of Tier~1 in the current ego frame, \emph{never} embeddings of future frames (which contain the ego's own motion and would teach an echo); an action-shuffled control whose failure voids Tier~2 checks this.
Tier~2 runs only for candidate families with counterfactual supervision, only on the survivors, and only when a \emph{learned} benefit gate fires (target $\le30\%$ of ticks). The disagreement of its two heads is a \emph{veto} on out-of-vocabulary candidates, not a benefit trigger: it detected them with an AUROC of $0.88$ (recall $64\%$ at $5\%$ false positives), while need-times-trust and margin gates were no better than random as triggers, and a learned gate reached $65\%$ of the gain at $30\%$ of the budget.

\paragraph{Composition and losses.}
\begin{equation}
\text{score}(a)=\tau\langle\mathbf u(s),\mathbf v(a)\rangle+\log\text{priors}-\beta_1J_1(s,a)-\beta_2J_2(s,a),\qquad J_1,J_2\ge0,
\end{equation}
so each imagination term adds a non-negative cost to a candidate (which can still change the argmin, since a candidate the base score ranked lower may be penalised less), and each term carries a zero-weight ablation. The imagination losses are $\mathcal L^{\text{img}}=\mathcal L^{\text{exo}}+\mathcal L^{\text{out}}+\lambda_{\text{inv}}\mathcal L^{\text{inv}}+\mathcal L^{\text{cons}}$: a Huber loss on descriptors plus cross-entropy on the mode plus a quantile loss on onset; an auxiliary outcome InfoNCE against the exogenous descriptors at $k\in\{5,10,20,40\}$ ticks, with per-state negatives where counterfactuals exist; an inverse-dynamics term of weight at most $0.1$ that keeps the action recoverable from state and future; and a parent--child hinge that keeps child outcomes consistent with their parent (strategic outcomes cannot be labelled without a map, so no strategic outcome quality is claimed).
\emph{No gradient flows from any imagination head into the frozen backbone}, and imagination parameters are trained separately from the selector's, since a world model that scores its own planner's fan was measured to help by only $+0.21$\evI.

\paragraph{Guards, each a measured failure.}
(G1)~An analytic feasibility and comfort mask is applied \emph{after} every learned term: imagination alone picked an infeasible over-braking candidate in $89.8\%$ of toy ticks and the mask brought it to $0\%$. (G2)~A counterfactual-coverage audit: a family without counterfactual supervision is refused by Tier 2. (G3)~A shuffled-cue and shuffled-$v_0$ control, with $R^2$ reported on the \emph{exogenous residual} only, because predicting the lead's speed from the ego state alone looks like imagination (echo). (G4)~Disagreement is a veto, never a benefit gate. (G5)~A confidence-sign audit: adding an imagined-cost term flipped the correlation of confidence with regret from $-0.11$ to $+0.19$ in all three seeds, so the combined softmax is never exposed as confidence. (G6)~$\beta$ is selected on held-out \emph{training} episodes that carry counterfactual labels (never on validation data), or it silently returns zero. (G7)~Every hidden factor is marginalised. (G8)~Every term is reported per family with a paired episode-cluster interval and its $n$ (agent-track coverage is stated for training clips only), and a risk-sensitive floor and a model-free value estimate on the same rows are always shown, otherwise risk-sensitivity is credited to imagination. (G9)~The goal path stays information-disjoint from the situation classifier: the E-head reads the state only.

\paragraph{Cost, and the first experiment.}
The imagination overhead is $\approx38$\,MFLOP per tick amortised (Tier 2 dominates), about $0.005\%$ of the smallest prefill (\evE); calling a learned world model for every candidate, or iterating rollouts over the whole set, would be $3$--$10\%$ of prefill and is not admitted. Efficiency is not the constraint; identifiability is.
The decisive first experiment is a probe costing $\approx0.1$\,GPU-hours on cached embeddings (E-I0): ridge and two-layer probes from the state to the lead gap now, the closing speed, the gap at $+2$\,s and the time to brake on the validation windows, against a $v_0$-only baseline and a shuffled-cue control. If the frozen state does not carry metric gap and closing speed (the programme's earlier frozen encoders did not), Tier 1 is inert and only Tier 0 ships; the toy's whole gap between $3.00$ and $0.46$ regret is exactly this readout error.

\paragraph{What could go wrong, and how the protocol sees it.}
(i)~Tier 1 may be inert because the backbone features lack range (E-I0, H-HC8). (ii)~The programme's earlier attempt to train an agent detector jointly with the trunk failed a two-seed capacity gate at $17$\,M parameters and the deranged-join control reproduced the whole degradation\evI; here the heads sit on a frozen backbone and do not compete with a trunk, but that is a hypothesis. (iii)~The non-reactive assumption fails at exactly the interactions that matter (yielding, merging); the one-sided form limits the damage and the shuffled controls (NC7) measure it. (iv)~In the toy the factorised arm \emph{lost} to the risk-sensitive floor when the cue was clean (regret $3.00$ against $2.17$ at a signal-to-noise ratio of $3$) because of visible-state readout error, the same failure as an earlier programme stream.
